\chapter{Interpreter Design}\label{ch:interpreter-design}

An interpreter must decide what a program means before any machine code exists.
It reads a representation of the program and carries out each operation directly.
That directness makes interpreters the fastest route to a working language and the slowest route to a fast program.
This chapter shows how to build one that is correct, predictable, and ready to support later optimization.
We treat the interpreter as an executable definition of the language semantics.
Every shortcut taken here becomes a commitment that later passes must respect.

\section{What an Interpreter Promises}\label{sec:ch30-promise}

A \term{interpreter} executes a program representation step by step without first translating it to native code.
The representation may be source text, an abstract syntax tree, or a compact bytecode.
The interpreter keeps explicit state for control, data, and environments.
Each step reads the current program point, performs the denoted action, and advances.
The observable result must match the language definition for every supported construct.
That promise covers values, effects, errors, and termination behavior.
A \term{tree-walking interpreter} traverses syntax nodes directly.
A \term{bytecode interpreter} dispatches on encoded instructions in a loop.
Both satisfy the same contract when they implement the same language.
They differ in startup cost, per-step cost, and preparation cost.
Consider a statement that adds two variables and stores the result.
A tree walker finds the assignment node, evaluates the right subtree, and updates the environment.
A bytecode interpreter fetches two loads, an add, and a store from an array.
The tree walker pays for pointer chasing and dispatch on node kinds.
The bytecode interpreter pays for decoding operands and threading control through one loop.
Neither cost matters until programs run long enough for dispatch overhead to dominate.
The choice therefore depends on how often each program runs and how quickly it must start.

\section{Structure of a Tree-Walking Interpreter}\label{sec:ch30-tree-walking}

A tree walker gives each syntactic form a node with an evaluation rule.
Evaluation takes a node and an \term{environment} that maps names to values or locations.
It returns a value, possibly after updating a store.
Recursive calls handle subexpressions in a defined order.
The order matters as soon as expressions can fail, diverge, or perform effects.
Figure~\ref{fig:ch30-eval-flow} shows how an evaluation request flows from a statement node through expression nodes to environment operations and back.
The diagram keeps control recursion separate from data lookup.
That separation simplifies later changes to scoping, assignment, and error reporting.
Environments deserve care because most interpreter bugs live there.
A flat map from names to values suffices for straight-line examples.
Nested scopes, functions, and closures need a chain of frames or a persistent mapping.
Assignment must update the frame where the name was bound, not the innermost frame.
A common mistake copies the environment on every call and loses updates to captured variables.
Another mistake shares one mutable map across unrelated activations and leaks bindings between calls.
We avoid both with an explicit rule.
Name resolution finds the nearest binding frame for a name.
Assignment writes to that frame.
Function creation captures the current frame chain for later calls.
This rule is short enough to test directly.
An interpreter loop also needs disciplined errors and termination control.
Division by zero, undefined names, and arity mismatches should name the offending construct and its source location.
The evaluator threads locations through every node so failures point at user code, not interpreter internals.
A recursion bound or an explicit value stack lets runaway programs fail cleanly.
These mechanisms are not decorations.
They define the observable behavior that compiled code must later reproduce.

\begin{figure}[htbp]
\centering
\begin{tikzpicture}[node distance=1.1cm and 1.4cm]
\node[basic block] (stmt) {assign node\\ \code{x = a + b}};
\node[basic block, below=of stmt] (add) {add node\\ \code{a + b}};
\node[basic block, below left=of add] (vara) {lookup\\ \code{a}};
\node[basic block, below right=of add] (varb) {lookup\\ \code{b}};
\node[basic block, below=of add, yshift=-1.6cm] (env) {environment\\ find, then update};
\draw[cfg edge] (stmt) -- (add);
\draw[cfg edge] (add) -- (vara);
\draw[cfg edge] (add) -- (varb);
\draw[cfg edge] (vara) -- (env);
\draw[cfg edge] (varb) -- (env);
\end{tikzpicture}
\caption{Evaluation of one assignment flows from the statement node to operand lookups and then to a single environment update, which keeps resolution and mutation in one place.}
\label{fig:ch30-eval-flow}
\end{figure}

\section{Dispatch: the Cost Center}\label{sec:ch30-dispatch}

Every interpreter spends most of its time deciding what to do next.
That decision is called \term{dispatch}.
A tree walker dispatches on node type at each recursive call.
A bytecode loop dispatches on opcode at each iteration.
Table~\ref{tab:ch30-dispatch} compares dispatch strategies by portability, effort, and branch behavior.
The table reports qualitative trade-offs rather than measured times.
Actual speed depends on the host compiler, the instruction set, and the program mix.
A switch inside a loop is the most portable choice.
Each iteration fetches an opcode and selects a case.
Computed jumps and threaded code reduce branch pressure but need host-specific support.
Direct call threading turns each handler into a small function and chains calls.
Subroutine threading emits native calls to handlers for straight-line sequences.
The progression trades portability for fewer unpredictable branches.
All strategies preserve the same semantics when handlers are unchanged.
We therefore fix handlers first and change dispatch only after profiling.
Algorithm~\ref{alg:ch30-bytecode-loop} shows the portable switch-dispatch loop reused by later chapters.
It assumes programs are arrays of records with an opcode and up to two operands.
It requires every opcode to have a handler and each handler to advance the program counter exactly once.
The loop checks errors, fuel exhaustion, and explicit halts on every iteration.
These checks look expensive in source form.
In practice they compile to a few predictable tests around handler bodies.
More importantly, they give the interpreter a defined stopping behavior that tests can observe.

\begin{table}[htbp]
\centering
\caption{Dispatch strategies trade portability against branch cost, so handlers should be fixed before dispatch is tuned.}
\label{tab:ch30-dispatch}
\begin{tabular}{llll}
\toprule
Strategy & Portability & Effort & Branch behavior \\
\midrule
Tree switch & High & Low & One branch per node type \\
Loop switch & High & Low & One branch per opcode \\
Computed jump & Medium & Medium & Indirect branch per opcode \\
Direct threading & Low & Medium & Chained handler jumps \\
Call threading & Medium & High & Native calls to handlers \\
\bottomrule
\end{tabular}
\end{table}

\begin{algorithm}[htbp]
\caption{Bytecode dispatch loop with explicit errors and fuel}\label{alg:ch30-bytecode-loop}
\begin{pseudocode}
\Require decoded program $P$ of $n$ instructions, state $S$, integer fuel $F\ge0$, terminating handlers
\Ensure halts with result, error, or out-of-fuel
\State $pc \gets 0$
\State $fuel \gets F$
\While{$pc \neq n$}
\If{$pc < 0$ or $pc > n$}
\State \Return invalid program counter
\EndIf
\If{$fuel = 0$}
\State \Return out-of-fuel
\EndIf
\State $fuel \gets fuel - 1$
\State $op \gets P[pc].opcode$
\If{$op$ is undefined}
\State \Return error at $pc$
\EndIf
\State execute checked handler on $S$; return any reported error immediately
\State accept its next $pc$ only if it is a valid target or the permitted end $n$
\EndWhile
\State \Return result in $S$
\end{pseudocode}
\end{algorithm}

This model treats $pc=n$ as successful fall-through; a format requiring an explicit halt must reject that transition instead.
Handlers check stack arity, operand kinds, indices, and arithmetic errors before publishing their next state.
Fuel bounds dispatches, not elapsed time: a handler that blocks forever or performs unbounded work defeats a wall-clock termination guarantee.

\section{Worked Example: Tracing One Assignment}\label{sec:ch30-worked}

We now trace the assignment \code{x = a + b} under concrete bindings.
This is a \term{hand-computed trace}, not a machine run.
Every value below follows from the previous step by the stated rules.
Suppose the environment maps $a$ to $6$ and $b$ to $9$.
Suppose $x$ is already bound in the current frame.
The tree walker visits the assignment node first.
It evaluates the right subtree before touching $x$.
It visits the addition node and evaluates $a$ to $6$.
It evaluates $b$ to $9$.
It adds the two values to get $15$.
It resolves $x$ to the current frame and stores $15$ there.
The final environment maps $a$ to $6$, $b$ to $9$, and $x$ to $15$.
The bytecode version performs the same steps in linear form.
The listing below uses abstract loads, an add, and a store.
Each instruction carries its operand inline so no decoding subtlety arises.
\begin{irlisting}
; x = a + b with a = 6 and b = 9
load a        ; stack: [6]
load b        ; stack: [6, 9]
add           ; stack: [15]
store x       ; env(x) = 15
\end{irlisting}
We step the dispatch loop of Algorithm~\ref{alg:ch30-bytecode-loop} with fuel $F = 10$.
The initial stack is empty and $pc = 0$.
After the first load, the stack holds $[6]$ and $pc = 1$.
After the second load, the stack holds $[6, 9]$ and $pc = 2$.
The add pops $9$ and $6$, pushes $15$, and sets $pc = 3$.
The store pops $15$, writes it to the frame for $x$, and sets $pc = 4$.
The loop exits because $pc$ equals the program length.
Fuel remaining is $6$, which confirms the trace consumed exactly four steps.
Both interpreters agree on the final store.
That agreement is the property tests should check.
If a later optimization changes dispatch, encoding, or evaluation order, this trace still defines the expected outcome.

\section{Values, Frames, and Error Paths}\label{sec:ch30-values}

A first interpreter should represent values with a tagged union or a small class hierarchy.
Integers, booleans, strings, and references need distinct tags even in a tiny language.
Operations check tags before computing and report a type error on mismatch.
That check belongs in the handler, not in dispatch.
Frames hold bindings, a return address, and a saved stack base.
Calls push a frame and reset the operand window.
Returns restore both and deliver one result value.
A frequent defect confuses the operand stack with the frame stack.
Keeping them as separate structures prevents a deep expression from corrupting return addresses.
Error paths need the same attention as value paths.
\begin{ccnote}[Testing errors]
Every error the definition allows needs a test that triggers it.
A missing test usually means the handler takes an undefined path instead of reporting cleanly.
\end{ccnote}
An implementation should test undefined names, arity mismatches, division guards, and stack underflow independently.
Each test asserts the diagnostic and the program point, not only the failure.
The \enquote{clean failure} rule keeps the suite deterministic.
No test may hang, crash the harness, or depend on host-specific output.
Fuel and depth limits enforce that rule mechanically.

\section{Implementation Notes and Complexity}\label{sec:ch30-implementation}

Represent program points as node pointers or integer indices, never as raw host addresses.
Represent environments as linked frames with hash maps for bindings.
These choices make inspection, printing, and testing straightforward.
The asymptotic cost is easy to state.
Evaluating an expression tree with $N$ nodes costs $O(N)$ handler steps when each node does constant work.
A bytecode loop over $I$ executed instructions costs $O(I)$ dispatches plus handler work.
Environment lookup costs $O(D)$ where $D$ is the static nesting depth for a frame chain.
A hash map per frame keeps each hop near constant time for typical programs.
Memory use grows with the value stack, the environment chain, and cached source locations.
Production systems depart from this baseline in measured steps.
They add inline caches for name and method lookup.
They specialize hot handlers for common types.
They replace tree walking with bytecode to shorten dispatch.
Each change keeps handler semantics fixed and measures the dispatch saving separately.
A common pitfall fuses these steps prematurely.
Fused handlers with ad hoc type tests are harder to verify than a slow interpreter with exact semantics.
Another pitfall omits fuel or stack limits for simplicity.
Without a limit, a diverging test hangs the suite instead of failing cleanly.
The interpreter doubles as the reference semantics for later chapters.
Its bytecode form supports the virtual machine of the next chapter.
Its environment discipline supports closures and objects later.
Its error and fuel behavior supports the testing discipline of the final chapters.
% TODO(cite): need a primary reference for threaded dispatch measurements.

\section{Exercises}\label{sec:ch30-exercises}

\begin{enumerate}
\item Trace the program \code{x = 2 + 3 + 4} under left-to-right evaluation and list the stack after every step, then repeat under right-to-left evaluation and state which final store, if any, differs.\label{ex:ch30-eval-order}
\item Implement the dispatch loop of Algorithm~\ref{alg:ch30-bytecode-loop} for four opcodes, add an out-of-fuel test with $F = 2$ on a four-instruction program, and report the program counter and stack at failure.\label{ex:ch30-fuel-test}
\item Extend a tree walker with one new binary operator, define its error behavior on two concrete inputs, and write tests that distinguish a missing handler from a wrong handler result.\label{ex:ch30-new-operator}
\end{enumerate}

\section{Further Reading}\label{sec:ch30-reading}

Appel presents interpreter structure as part of a larger compiler pipeline~\cite{appel-1998-modern}.
Morgan discusses how interpreter and runtime organization constrains later optimization~\cite{morgan-1998-optimizing}.
Safonov treats language definition and trustworthy execution as engineering obligations~\cite{safonov-2010-trustworthy}.
The C rules for expression evaluation and sequencing are defined in the language standard draft~\cite{wg14-2011-c11}.
